% Part: lambda-calculus
% Chapter: lambda-definability
% Section: lambda-definable-recursive

\documentclass[../../../include/open-logic-section]{subfiles}

\begin{document}

\olfileid{lam}{ldf}{ldr}
\olsection{\usetoken{S}{lambda definable} વિધેયો પુનરાવર્તી છે}
\sourcecorrection{OLLAM-062}{સ્થિર અંગ્રેજી મૂળમાં આ
વિભાગની અધ્યાય-ઓળખ~\texttt{dfl} છે, જ્યારે
અધ્યાયની ઓળખ~\texttt{ldf} અને અન્ય વિભાગોની
ઓળખ પણ~\texttt{ldf} છે. આંતરિક સંદર્ભો માટે
અહીં~\texttt{ldf} કર્યું છે.}

બધા આંશિક પુનરાવર્તી વિધેયો !!{lambda definable}
છે એટલું જ નહીં, તેનો વિપરીત દાવો પણ સાચો છે.
એટલે બધા !!{lambda definable} વિધેયો આંશિક
પુનરાવર્તી છે.

\begin{thm}
  \ollabel{thm:lambda-computable} આંશિક વિધેય~$f$
  !!{lambda definable} હોય તો તે આંશિક પુનરાવર્તી છે.
\end{thm}

\begin{proof}
  પુરાવાની માત્ર રૂપરેખા આપીશું. પહેલા આપણે
  $\lambd$-પદોનું અંકગણિતીકરણ કરીએ છીએ: અનુક્રમોને
  અભાજ્ય સંખ્યાઓની ઘાતોના ગુણાકાર તરીકે સંકેતિત કરવાની પ્રચલિત રીતથી
  $\lambd$-પદોને વ્યવસ્થિત રીતે ગોડેલ સંખ્યાઓ આપીએ
  છીએ. પછી આંશિક પુનરાવર્તી વિધેય
  $\fn{normalize}(t)$ વ્યાખ્યાયિત કરીએ છીએ. તે
  લૅમ્બડા પદની ગોડેલ સંખ્યા~$t$ને આર્ગ્યુમેન્ટ
  તરીકે લે છે; પદનું સામાન્ય સ્વરૂપ હોય તો તેની
  ગોડેલ સંખ્યા પરત કરે છે અને ન હોય તો
  અવ્યાખ્યાયિત રહે છે. ત્યાર પછી $\fn{toChurch}$
  અને $\fn{fromChurch}$ એવાં બે આંશિક પુનરાવર્તી
  વિધેયો વ્યાખ્યાયિત કરીએ છીએ, જે પ્રાકૃતિક સંખ્યાઓને
  અનુરૂપ ચર્ચ અંકોની ગોડેલ સંખ્યાઓમાં અને તેમાંથી
  પાછી ફેરવે છે.

  આ પુનરાવર્તી વિધેયો વડે $f$ને આંશિક પુનરાવર્તી
  વિધેય તરીકે વ્યાખ્યાયિત કરી શકીએ છીએ. કોઈ
  $\lambd$-પદ~$F$ વિધેય~$f$ને !!{lambda define}
  કરે છે. $f(n_1, \dots, n_k)$ ગણવા પહેલા
  $\fn{toChurch}(n_i)$ વડે અનુરૂપ ચર્ચ અંકોની
  ગોડેલ સંખ્યાઓ મેળવો. તેમને $\Gn{F}$ સાથે જોડીને
  $F \num{n_1}\dots\num{n_k}$ પદની ગોડેલ સંખ્યા
  બનાવો. હવે આ ગોડેલ સંખ્યા પર $\fn{normalize}$
  લાગુ કરો. $f(n_1, \dots, n_k)$ વ્યાખ્યાયિત હોય
  તો $F \num{n_1}\dots\num{n_k}$નું સામાન્ય સ્વરૂપ
  હોય છે (અને તે ચર્ચ અંક જ હોવો જોઈએ); અન્યથા
  તેનું સામાન્ય સ્વરૂપ હોતું નથી (એટલે
  \[\fn{normalize}(\Gn{F\num{n_1}\dots\num{n_k}})\]
  અવ્યાખ્યાયિત છે). છેવટે સામાન્ય સ્વરૂપે આવેલા
  પદની ગોડેલ સંખ્યા પર $\fn{fromChurch}$ લાગુ કરો.
\end{proof}

\end{document}
